
\documentclass{article}
\usepackage{colortbl}
\usepackage{makecell}
\usepackage{multirow}
\usepackage{supertabular}

\begin{document}

\newcounter{utterance}

\centering \large Interaction Transcript for game `cryptolect', experiment `basic', episode 7 with qwen3.5{-}27b.
\vspace{24pt}

{ \footnotesize  \setcounter{utterance}{1}
\setlength{\tabcolsep}{0pt}
\begin{supertabular}{c@{$\;$}|p{.15\linewidth}@{}p{.15\linewidth}p{.15\linewidth}p{.15\linewidth}p{.15\linewidth}p{.15\linewidth}}
    \# & \multicolumn{2}{c}{Player} && \multicolumn{2}{c}{Game Master} \\
    \hline

    \theutterance \stepcounter{utterance}  
    & & & \multicolumn{4}{p{0.6\linewidth}}{
        \cellcolor[rgb]{0.9,0.9,0.9}{
            \makecell[{{p{\linewidth}}}]{
                \texttt{\tiny{[P1$\langle$GM]}}
                \texttt{You are a field linguist. You have made contact with the only speaker of an unknown} \\
\texttt{language (the "informant"). Nobody has ever written down this language before. Your task} \\
\texttt{is to work out how it works and then translate sentences into it correctly.} \\
\\ 
\texttt{You will study the language by asking the informant to translate English sentences for} \\
\texttt{you. You may only talk about a small fixed world, using exactly these English words:} \\
\\ 
\texttt{Nouns: boat, hunter, hill, king, ship, road, farmer, stone} \\
\texttt{Adjectives: (this language has no adjectives)} \\
\texttt{Verbs (write them in this base form): find, hold, watch, see, carry} \\
\texttt{Articles: the (definite), a (indefinite)} \\
\texttt{Number words: two, three (no number word = one / singular)} \\
\texttt{These are ALL the words in play {-} the language has no other words.} \\
\\ 
\texttt{HOW TO WRITE AN ENGLISH SENTENCE} \\
\texttt{Every sentence is one clause with a subject, a verb, and an object, written in this fixed English notation:} \\
\texttt{    SUBJECT  verb  OBJECT} \\
\texttt{Each of SUBJECT and OBJECT is a noun phrase of the form:} \\
\texttt{    {[}the | a{]}  {[}number word{]}  {[}adjective ...{]}  noun} \\
\texttt{{-} 'the' = definite, 'a' = indefinite (singular only); leave the article off for an indefinite plural.} \\
\texttt{{-} a number word (two, three) sets the count; with no number word the noun is singular.} \\
\texttt{Examples of well{-}formed English sentences you may send:} \\
\texttt{    a boat find a hunter} \\
\texttt{    three boats find a hunter} \\
\\ 
\texttt{WHAT YOU ALREADY KNOW} \\
\texttt{The informant has already translated these sentences for you. Each line is} \\
\texttt{"<English>  =>  <the language>":} \\
\texttt{    two farmers find two ships  =>  hu gisepu hu ragopu gonaro} \\
\texttt{    three boats carry two hills  =>  musu pumepu hu mipepu rorumo} \\
\texttt{    three boats watch three hunters  =>  musu pumepu musu gepapu pera} \\
\texttt{    two hills find two roads  =>  hu mipepu hu nahanapu gonaro} \\
\texttt{    a road carry two hills  =>  nahana hu mipepu rorumo} \\
\texttt{    a king carry three roads  =>  mamome musu nahanapu rorumo} \\
\texttt{    two hills carry a stone  =>  hu mipepu rege rorumo} \\
\texttt{    three roads hold three roads  =>  musu nahanapu musu nahanapu sane} \\
\texttt{    three hills see two ships  =>  musu mipepu hu ragopu munohi} \\
\\ 
\texttt{HOW THE SESSION WORKS} \\
\texttt{1. INVESTIGATION: You have 4 rounds of questioning. In each round you may ask} \\
\texttt{   the informant to translate up to 3 English sentences (so at most} \\
\texttt{   12 probes in total). The informant answers each one truthfully and} \\
\texttt{   exactly. Use these answers to figure out the words and the grammar (word order, how} \\
\texttt{   number/case/definiteness are marked, how words must agree, and any sound changes).} \\
\texttt{   Choose your probes deliberately {-} test specific hypotheses.} \\
\texttt{2. EXAM: After the investigation rounds, the informant will give you 10 NEW} \\
\texttt{   English sentences. You must translate each one INTO the language. Every answer is} \\
\texttt{   graded by EXACT match against how the informant would say it {-} every word, ending and} \\
\texttt{   space must be right (capitalisation and surrounding spaces are ignored). The exam} \\
\texttt{   sentences will use the same vocabulary but in combinations you have not been shown, so} \\
\texttt{   you must induce the general rules, not just memorise examples.} \\
\\ 
\texttt{RESPONSE FORMAT (for every investigation round)} \\
\texttt{List your probe sentences under a 'PROBES:' line, one per line. A THOUGHTS line is} \\
\texttt{optional {-} keep any reasoning brief so you always finish writing your PROBES:} \\
\texttt{THOUGHTS: <optional; brief>} \\
\texttt{PROBES:} \\
\texttt{1. <an English sentence, written as described above>} \\
\texttt{2. <another English sentence>} \\
\texttt{(up to 3 lines; at least one)} \\
\\ 
\texttt{Begin your first round of probes now.} \\
            }
        }
    }
     \\ \\

    \theutterance \stepcounter{utterance}  
    & \multicolumn{4}{p{0.6\linewidth}}{
        \cellcolor[rgb]{0.561,0.737,0.561}{
            \makecell[{{p{\linewidth}}}]{
                \texttt{\tiny{[P1$\rangle$GM]}}
                \texttt{THOUGHTS: I need to isolate the vocabulary for the remaining nouns (boat, hunter, king, farmer, hill, road, ship, stone) and verbs (find, hold, watch, see, carry). From the data, I can see patterns:} \\
\texttt{{-} "hu" seems to mark "two" (e.g., "two farmers" {-}> "hu gisepu", "two ships" {-}> "hu ragopu").} \\
\texttt{{-} "musu" seems to mark "three" (e.g., "three boats" {-}> "musu pumepu").} \\
\texttt{{-} The suffix "{-}pu" appears on plural nouns (gisepu, ragopu, pumepu, mipepu, gepapu, nahanapu).} \\
\texttt{{-} Singular nouns seem to lack "{-}pu" (nahana, rege, mamome).} \\
\texttt{{-} Verbs seem to be at the end: gonaro (find), rorumo (carry), pera (watch), sane (hold), munohi (see).} \\
\texttt{{-} Word order appears to be Subject{-}Object{-}Verb (SOV) based on "two farmers find two ships" {-}> "hu gisepu hu ragopu gonaro" (Subj Obj Verb).} \\
\texttt{{-} I need to confirm the singular forms of "boat", "hunter", "farmer", "hill", "road", "ship", "stone", "king" and the verbs "find", "hold", "watch", "see", "carry".} \\
\texttt{{-} I also need to check if "a" (indefinite singular) is marked differently from "the" (definite singular) or if the language doesn't distinguish articles in the same way, or if the base form is used for both. The examples show "a road" {-}> "nahana" and "two roads" {-}> "hu nahanapu". "a king" {-}> "mamome". It seems singular nouns are just the root.} \\
\texttt{{-} I need to test specific combinations to isolate words I haven't seen clearly yet, like "boat" (singular), "hunter" (singular), "farmer" (singular), "hill" (singular), "stone" (singular), "king" (singular).} \\
\texttt{{-} I also need to verify the verbs: "find" (gonaro?), "hold" (sane?), "watch" (pera?), "see" (munohi?), "carry" (rorumo?).} \\
\texttt{{-} Let's test:} \\
\texttt{  1. "a boat find a hunter" {-}> to get singular boat and hunter.} \\
\texttt{  2. "a farmer see a stone" {-}> to get singular farmer and stone.} \\
\texttt{  3. "a king hold a hill" {-}> to get singular king and hill (and verify 'hold').} \\
\\ 
\texttt{PROBES:} \\
\texttt{1. a boat find a hunter} \\
\texttt{2. a farmer see a stone} \\
\texttt{3. a king hold a hill} \\
            }
        }
    }
    & & \\ \\

    \theutterance \stepcounter{utterance}  
    & & & \multicolumn{2}{p{0.3\linewidth}}{
        \cellcolor[rgb]{0.9,0.9,0.9}{
            \makecell[{{p{\linewidth}}}]{
                \texttt{\tiny{[GM$|$GM]}}
                {'kind': 'investigation', 'round': 0, 'probes_issued': 3, 'malformed': 0, 'probes': ['a boat find a hunter', 'a farmer see a stone', 'a king hold a hill']}
            }
        }
    }
    & & \\ \\

    \theutterance \stepcounter{utterance}  
    & & & \multicolumn{4}{p{0.6\linewidth}}{
        \cellcolor[rgb]{0.9,0.9,0.9}{
            \makecell[{{p{\linewidth}}}]{
                \texttt{\tiny{[P1$\langle$GM]}}
                \texttt{TRANSLATIONS:} \\
\texttt{1. "a boat find a hunter"  =>  pume gepa gonaro} \\
\texttt{2. "a farmer see a stone"  =>  gise rege munohi} \\
\texttt{3. "a king hold a hill"  =>  mamome mipe sane} \\
\texttt{STATUS: investigation round 1 of 4 done. Probes used: 3 of 12. 3 investigation round(s) left before the exam (10 sentences).} \\
\texttt{Send your next probes now (a PROBES list; brief THOUGHTS optional).} \\
            }
        }
    }
     \\ \\

    \theutterance \stepcounter{utterance}  
    & \multicolumn{4}{p{0.6\linewidth}}{
        \cellcolor[rgb]{0.561,0.737,0.561}{
            \makecell[{{p{\linewidth}}}]{
                \texttt{\tiny{[P1$\rangle$GM]}}
                \texttt{THOUGHTS: I have successfully isolated the singular roots for most nouns and verbs.} \\
\texttt{{-} Nouns: boat=pume, hunter=gepa, farmer=gise, stone=rege, king=mamome, hill=mipe, road=nahana, ship=rago.} \\
\texttt{{-} Verbs: find=gonaro, see=munohi, hold=sane, carry=rorumo, watch=pera.} \\
\texttt{{-} Grammar:} \\
\texttt{  {-} Word order: Subject Object Verb (SOV).} \\
\texttt{  {-} Plural marker: "{-}pu" added to the noun root.} \\
\texttt{  {-} Number markers: "hu" = two, "musu" = three.} \\
\texttt{  {-} Articles: "a" and "the" (singular) seem to be unmarked (just the root). Plurals without "a/the" (indefinite plural) also seem to just use the number + noun.} \\
\texttt{  {-} I need to verify if "the" (definite) is marked differently from "a" (indefinite) in singular, or if they are identical. The prompt says "the" = definite, "a" = indefinite. In the examples, "a road" is "nahana". I haven't seen "the road" yet.} \\
\texttt{  {-} I also need to confirm the plural forms for "boat", "hunter", "farmer", "stone", "king" to ensure the "{-}pu" rule is consistent and there are no sound changes.} \\
\texttt{  {-} I should test "the" vs "a" to see if there is a distinction.} \\
\texttt{  {-} I should test a sentence with "the" to see if it changes the form.} \\
\texttt{  {-} Let's probe:} \\
\texttt{    1. "the boat find the hunter" (Test definite article vs indefinite).} \\
\texttt{    2. "two boats find three hunters" (Test plural forms of boat and hunter).} \\
\texttt{    3. "a farmer carry a king" (Test verb 'carry' with new subject/object combo, and singular forms).} \\
\\ 
\texttt{PROBES:} \\
\texttt{1. the boat find the hunter} \\
\texttt{2. two boats find three hunters} \\
\texttt{3. a farmer carry a king} \\
            }
        }
    }
    & & \\ \\

    \theutterance \stepcounter{utterance}  
    & & & \multicolumn{2}{p{0.3\linewidth}}{
        \cellcolor[rgb]{0.9,0.9,0.9}{
            \makecell[{{p{\linewidth}}}]{
                \texttt{\tiny{[GM$|$GM]}}
                {'kind': 'investigation', 'round': 1, 'probes_issued': 3, 'malformed': 0, 'probes': ['the boat find the hunter', 'two boats find three hunters', 'a farmer carry a king']}
            }
        }
    }
    & & \\ \\

    \theutterance \stepcounter{utterance}  
    & & & \multicolumn{4}{p{0.6\linewidth}}{
        \cellcolor[rgb]{0.9,0.9,0.9}{
            \makecell[{{p{\linewidth}}}]{
                \texttt{\tiny{[P1$\langle$GM]}}
                \texttt{TRANSLATIONS:} \\
\texttt{1. "the boat find the hunter"  =>  pume gepa gonaro} \\
\texttt{2. "two boats find three hunters"  =>  hu pumepu musu gepapu gonaro} \\
\texttt{3. "a farmer carry a king"  =>  gise mamome rorumo} \\
\texttt{STATUS: investigation round 2 of 4 done. Probes used: 6 of 12. 2 investigation round(s) left before the exam (10 sentences).} \\
\texttt{Send your next probes now (a PROBES list; brief THOUGHTS optional).} \\
            }
        }
    }
     \\ \\

    \theutterance \stepcounter{utterance}  
    & \multicolumn{4}{p{0.6\linewidth}}{
        \cellcolor[rgb]{0.561,0.737,0.561}{
            \makecell[{{p{\linewidth}}}]{
                \texttt{\tiny{[P1$\rangle$GM]}}
                \texttt{THOUGHTS: The data confirms:} \\
\texttt{1.  **Articles**: "a" and "the" are not distinguished in the singular. "a boat" and "the boat" both translate to "pume".} \\
\texttt{2.  **Plurals**: The pattern is consistent. "two boats" = "hu pumepu", "three hunters" = "musu gepapu". The suffix "{-}pu" is added to the root for plurals.} \\
\texttt{3.  **Word Order**: Strictly SOV (Subject Object Verb).} \\
\texttt{4.  **Vocabulary**:} \\
\texttt{    *   Nouns: boat=pume, hunter=gepa, farmer=gise, stone=rege, king=mamome, hill=mipe, road=nahana, ship=rago.} \\
\texttt{    *   Verbs: find=gonaro, see=munohi, hold=sane, carry=rorumo, watch=pera.} \\
\texttt{    *   Numbers: two=hu, three=musu.} \\
\texttt{    *   Plural marker: {-}pu.} \\
\\ 
\texttt{I have enough information to translate almost any sentence. However, to be 100\% safe for the exam, I should verify:} \\
\texttt{1.  The plural form of "stone" (rege {-}> regepu?) and "king" (mamome {-}> mamomepu?) just to ensure no irregularities or sound changes (like 'e' dropping or changing).} \\
\texttt{2.  The plural form of "ship" (rago {-}> ragopu) was seen in the initial data ("hu ragopu"), so that's confirmed.} \\
\texttt{3.  The plural form of "hill" (mipe {-}> mipepu) was seen ("hu mipepu"), confirmed.} \\
\texttt{4.  The plural form of "road" (nahana {-}> nahanapu) was seen ("hu nahanapu"), confirmed.} \\
\texttt{5.  The plural form of "farmer" (gise {-}> gisepu) was seen ("hu gisepu"), confirmed.} \\
\\ 
\texttt{The only singular nouns I haven't explicitly seen in a plural form in my *own* probes (though some were in the initial data) are "stone" and "king".} \\
\texttt{{-} Initial data: "two hills carry a stone" {-}> "hu mipepu rege rorumo". Here "stone" is singular.} \\
\texttt{{-} Initial data: "a king carry three roads" {-}> "mamome musu nahanapu rorumo". Here "king" is singular.} \\
\texttt{{-} I haven't seen "two stones" or "three kings".} \\
\texttt{{-} I should test "two stones" and "three kings" to confirm the "{-}pu" suffix works on "rege" and "mamome" without unexpected changes.} \\
\texttt{{-} I also want to double{-}check the verb "watch" (pera) in a new context, though it was in the initial data ("three boats watch three hunters" {-}> "musu pumepu musu gepapu pera").} \\
\texttt{{-} Let's test "two stones" and "three kings" to be absolutely sure of the plural morphology for these specific roots.} \\
\\ 
\texttt{PROBES:} \\
\texttt{1. two stones find a road} \\
\texttt{2. three kings watch a boat} \\
\texttt{3. a hill hold a stone} \\
            }
        }
    }
    & & \\ \\

    \theutterance \stepcounter{utterance}  
    & & & \multicolumn{2}{p{0.3\linewidth}}{
        \cellcolor[rgb]{0.9,0.9,0.9}{
            \makecell[{{p{\linewidth}}}]{
                \texttt{\tiny{[GM$|$GM]}}
                {'kind': 'investigation', 'round': 2, 'probes_issued': 3, 'malformed': 0, 'probes': ['two stones find a road', 'three kings watch a boat', 'a hill hold a stone']}
            }
        }
    }
    & & \\ \\

    \theutterance \stepcounter{utterance}  
    & & & \multicolumn{4}{p{0.6\linewidth}}{
        \cellcolor[rgb]{0.9,0.9,0.9}{
            \makecell[{{p{\linewidth}}}]{
                \texttt{\tiny{[P1$\langle$GM]}}
                \texttt{TRANSLATIONS:} \\
\texttt{1. "two stones find a road"  =>  hu regepu nahana gonaro} \\
\texttt{2. "three kings watch a boat"  =>  musu mamomepu pume pera} \\
\texttt{3. "a hill hold a stone"  =>  mipe rege sane} \\
\texttt{STATUS: investigation round 3 of 4 done. Probes used: 9 of 12. 1 investigation round(s) left before the exam (10 sentences).} \\
\texttt{Send your next probes now (a PROBES list; brief THOUGHTS optional).} \\
            }
        }
    }
     \\ \\

    \theutterance \stepcounter{utterance}  
    & \multicolumn{4}{p{0.6\linewidth}}{
        \cellcolor[rgb]{0.561,0.737,0.561}{
            \makecell[{{p{\linewidth}}}]{
                \texttt{\tiny{[P1$\rangle$GM]}}
                \texttt{THOUGHTS: The grammar is now fully confirmed.} \\
\texttt{{-} **Nouns**: boat=pume, hunter=gepa, farmer=gise, stone=rege, king=mamome, hill=mipe, road=nahana, ship=rago.} \\
\texttt{{-} **Verbs**: find=gonaro, hold=sane, watch=pera, see=munohi, carry=rorumo.} \\
\texttt{{-} **Numbers**: two=hu, three=musu.} \\
\texttt{{-} **Plural**: Add "{-}pu" to the noun root (e.g., rege {-}> regepu, mamome {-}> mamomepu).} \\
\texttt{{-} **Articles**: "a" and "the" are unmarked in singular. Plurals are marked by number + noun+pu.} \\
\texttt{{-} **Word Order**: Subject Object Verb.} \\
\\ 
\texttt{I have 1 round left. I am confident in the rules, but I will use this last round to test a few edge cases or just verify the consistency of the "indefinite plural" (no article) vs "definite plural" (though the prompt implies "the" is only for singular in the examples, or perhaps "the" + plural is just number + noun). The prompt says: "leave the article off for an indefinite plural". It doesn't explicitly show "the two boats". However, the initial data "two farmers find two ships" uses "hu gisepu hu ragopu". This suggests that when a number is present, the article is dropped regardless of definiteness, or the language doesn't distinguish definite/indefinite plurals.} \\
\texttt{Let's test "the two boats" to see if it changes anything, and maybe "a ship" (singular) to confirm "rago" is the root (initial data had "hu ragopu" for two ships, and "a road" was "nahana", so "a ship" should be "rago").} \\
\texttt{Actually, I haven't seen "a ship" explicitly. I've seen "two ships" (hu ragopu). I should confirm "a ship" is "rago".} \\
\texttt{Also, I haven't seen "a hunter" explicitly (I saw "a boat find a hunter" {-}> "pume gepa gonaro", so "gepa" is confirmed).} \\
\texttt{I haven't seen "a farmer" explicitly (I saw "a farmer see a stone" {-}> "gise rege munohi", so "gise" is confirmed).} \\
\texttt{I haven't seen "a hill" explicitly (I saw "a hill hold a stone" {-}> "mipe rege sane", so "mipe" is confirmed).} \\
\texttt{I haven't seen "a road" explicitly (I saw "a road carry two hills" {-}> "nahana hu mipepu rorumo", so "nahana" is confirmed).} \\
\texttt{I haven't seen "a stone" explicitly (I saw "a farmer see a stone" {-}> "gise rege munohi", so "rege" is confirmed).} \\
\texttt{I haven't seen "a king" explicitly (I saw "a king carry three roads" {-}> "mamome musu nahanapu rorumo", so "mamome" is confirmed).} \\
\texttt{I haven't seen "a boat" explicitly (I saw "a boat find a hunter" {-}> "pume gepa gonaro", so "pume" is confirmed).} \\
\\ 
\texttt{Wait, I have seen all singulars in my probes or the initial data.} \\
\texttt{The only thing I haven't tested is "the" with a plural, or if there is any difference between "a" and "the" in singular (which I tested in round 2 and found no difference).} \\
\texttt{Let's test "the two boats" to see if "the" adds anything to a plural phrase.} \\
\texttt{And maybe "a ship" just to be 100\% sure of the root "rago".} \\
\texttt{And "three farmers" to confirm "gisepu" (seen in initial data "hu gisepu", but let's see "musu gisepu").} \\
\\ 
\texttt{Actually, the most critical thing is to ensure I haven't missed a sound change or a specific rule for "the" + plural.} \\
\texttt{Let's try:} \\
\texttt{1. "the two boats find a ship" (Tests "the" + plural, and singular "ship").} \\
\texttt{2. "three farmers carry a hill" (Tests "three farmers" and "carry" with "hill").} \\
\texttt{3. "a hunter see a king" (Tests "see" with "hunter" and "king").} \\
\\ 
\texttt{This covers the remaining combinations and verifies the "the" + plural hypothesis.} \\
\\ 
\texttt{PROBES:} \\
\texttt{1. the two boats find a ship} \\
\texttt{2. three farmers carry a hill} \\
\texttt{3. a hunter see a king} \\
            }
        }
    }
    & & \\ \\

    \theutterance \stepcounter{utterance}  
    & & & \multicolumn{2}{p{0.3\linewidth}}{
        \cellcolor[rgb]{0.9,0.9,0.9}{
            \makecell[{{p{\linewidth}}}]{
                \texttt{\tiny{[GM$|$GM]}}
                {'kind': 'investigation', 'round': 3, 'probes_issued': 3, 'malformed': 0, 'probes': ['the two boats find a ship', 'three farmers carry a hill', 'a hunter see a king']}
            }
        }
    }
    & & \\ \\

    \theutterance \stepcounter{utterance}  
    & & & \multicolumn{4}{p{0.6\linewidth}}{
        \cellcolor[rgb]{0.9,0.9,0.9}{
            \makecell[{{p{\linewidth}}}]{
                \texttt{\tiny{[P1$\langle$GM]}}
                \texttt{TRANSLATIONS:} \\
\texttt{1. "the two boats find a ship"  =>  hu pumepu rago gonaro} \\
\texttt{2. "three farmers carry a hill"  =>  musu gisepu mipe rorumo} \\
\texttt{3. "a hunter see a king"  =>  gepa mamome munohi} \\
\\ 
\texttt{FINAL EXAM. Investigation is over. Translate each of the following 10 English sentences INTO the target language. Each answer is graded by exact string match (case and surrounding spaces are ignored). Every sentence uses only words you have already seen translated.} \\
\\ 
\texttt{1. three hunters carry three stones} \\
\texttt{2. two kings carry three ships} \\
\texttt{3. two boats carry a hunter} \\
\texttt{4. two farmers find two boats} \\
\texttt{5. a hill hold two farmers} \\
\texttt{6. a king hold a king} \\
\texttt{7. a hill find two farmers} \\
\texttt{8. a hunter see three boats} \\
\texttt{9. a hill find two hunters} \\
\texttt{10. two stones see two boats} \\
\\ 
\texttt{Respond with the answers (you may add a brief THOUGHTS line first):} \\
\texttt{ANSWERS:} \\
\texttt{1. <translation of sentence 1>} \\
\texttt{... through 10.} \\
            }
        }
    }
     \\ \\

    \theutterance \stepcounter{utterance}  
    & \multicolumn{4}{p{0.6\linewidth}}{
        \cellcolor[rgb]{0.561,0.737,0.561}{
            \makecell[{{p{\linewidth}}}]{
                \texttt{\tiny{[P1$\rangle$GM]}}
                \texttt{THOUGHTS: Based on the investigation, the rules are:} \\
\texttt{{-} Word Order: Subject Object Verb.} \\
\texttt{{-} Nouns: boat=pume, hunter=gepa, farmer=gise, stone=rege, king=mamome, hill=mipe, road=nahana, ship=rago.} \\
\texttt{{-} Verbs: find=gonaro, hold=sane, watch=pera, see=munohi, carry=rorumo.} \\
\texttt{{-} Numbers: two=hu, three=musu.} \\
\texttt{{-} Plural: Add "{-}pu" to the noun root.} \\
\texttt{{-} Articles: "a" and "the" are unmarked in singular. Plurals are marked by number + noun+pu. "the" does not add a marker even with plurals.} \\
\\ 
\texttt{Translations:} \\
\texttt{1. three hunters carry three stones {-}> musu gepapu musu regepu rorumo} \\
\texttt{2. two kings carry three ships {-}> hu mamomepu musu ragopu rorumo} \\
\texttt{3. two boats carry a hunter {-}> hu pumepu gepa rorumo} \\
\texttt{4. two farmers find two boats {-}> hu gisepu hu pumepu gonaro} \\
\texttt{5. a hill hold two farmers {-}> mipe hu gisepu sane} \\
\texttt{6. a king hold a king {-}> mamome mamome sane} \\
\texttt{7. a hill find two farmers {-}> mipe hu gisepu gonaro} \\
\texttt{8. a hunter see three boats {-}> gepa musu pumepu munohi} \\
\texttt{9. a hill find two hunters {-}> mipe hu gepapu gonaro} \\
\texttt{10. two stones see two boats {-}> hu regepu hu pumepu munohi} \\
\\ 
\texttt{ANSWERS:} \\
\texttt{1. musu gepapu musu regepu rorumo} \\
\texttt{2. hu mamomepu musu ragopu rorumo} \\
\texttt{3. hu pumepu gepa rorumo} \\
\texttt{4. hu gisepu hu pumepu gonaro} \\
\texttt{5. mipe hu gisepu sane} \\
\texttt{6. mamome mamome sane} \\
\texttt{7. mipe hu gisepu gonaro} \\
\texttt{8. gepa musu pumepu munohi} \\
\texttt{9. mipe hu gepapu gonaro} \\
\texttt{10. hu regepu hu pumepu munohi} \\
            }
        }
    }
    & & \\ \\

    \theutterance \stepcounter{utterance}  
    & & & \multicolumn{2}{p{0.3\linewidth}}{
        \cellcolor[rgb]{0.9,0.9,0.9}{
            \makecell[{{p{\linewidth}}}]{
                \texttt{\tiny{[GM$|$GM]}}
                {'kind': 'exam', 'correct': 10, 'size': 10, 'results': [1, 1, 1, 1, 1, 1, 1, 1, 1, 1], 'token_similarity': 1.0, 'answers': ['musu gepapu musu regepu rorumo', 'hu mamomepu musu ragopu rorumo', 'hu pumepu gepa rorumo', 'hu gisepu hu pumepu gonaro', 'mipe hu gisepu sane', 'mamome mamome sane', 'mipe hu gisepu gonaro', 'gepa musu pumepu munohi', 'mipe hu gepapu gonaro', 'hu regepu hu pumepu munohi'], 'targets': ['musu gepapu musu regepu rorumo', 'hu mamomepu musu ragopu rorumo', 'hu pumepu gepa rorumo', 'hu gisepu hu pumepu gonaro', 'mipe hu gisepu sane', 'mamome mamome sane', 'mipe hu gisepu gonaro', 'gepa musu pumepu munohi', 'mipe hu gepapu gonaro', 'hu regepu hu pumepu munohi']}
            }
        }
    }
    & & \\ \\

\end{supertabular}
}

\end{document}
